\chapter{نتیجه‌گیری و جمع‌بندی}
\section{مرور دستاوردها}
در این پروژه، الگوریتم گرادیان کاهشی به عنوان یکی از روش‌های بنیادین بهینه‌سازی محدب معرفی و پیاده‌سازی شد. مهم‌ترین دستاوردها عبارتند از:
\begin{itemize}
	\item ارائه سه فرمول کلیدی: فرمول بازگشتی گرادیان کاهشی در قالب \lr{align}، تابع جریمه چندضابطه‌ای و فرمول نرم اقلیدسی.
	\item نوشتن شبه‌کد الگوریتم گرادیان کاهشی در محیط \lr{algorithm} و رسم فلوچارت مربوطه با استفاده از \lr{TikZ}.
	\item طراحی جدول مقایسه‌ای با رعایت جزئیات خواسته شده (خطوط ضخیم، رنگ‌آمیزی یک‌درمیان، تنظیم ابعاد سلول).
	\item ایجاد شکل با دو زیرشکل برای نمایش اثر نرخ یادگیری بر همگرایی.
	\item طراحی گراف تعاملی برای نشان دادن ارتباط بین فصول با قابلیت کلیک (در این فصل ارائه شده است).
	\item رعایت کامل اصول نگارشی (نقطه پایان جمله، نیم‌فاصله، استفاده از محیط ریاضی برای نمادها، پرهیز از کلمات انگلیسی بدون پانوشت).
\end{itemize}

\section{نتایج عددی شبیه‌سازی}
به منظور ارزیابی الگوریتم، تابع هدف \( f(x) = x^2 + 2x + 1 \) در نظر گرفته شد. با نقطه شروع \( x_0 = 5 \) و نرخ یادگیری \( \alpha = 0.1 \) و معیار توقف \( \varepsilon = 10^{-6} \)، الگوریتم پس از ۸۰ تکرار به نقطه بهینه \( x^* = -1.00 \) با مقدار تابع \( f(x^*) \approx 0 \) همگرا شد. در مقابل، با نرخ یادگیری \( \alpha = 0.9 \) الگوریتم دچار نوسان و همگرایی کندتری شد که در زیرشکل‌های \cref{fig:subfigures} به تصویر کشیده شده است.

\section{گراف تعاملی فصول}
در شکل \cref{fig:graph-chapters} گرافی طراحی شده که با کلیک روی هر نود به فصل مربوطه منتقل می‌شوید. این قابلیت با استفاده از بسته \lr{hyperref} و دستور \lr{\hyperlink} درون گره‌های \lr{TikZ} پیاده‌سازی شده است.
\begin{figure}[H]
	\centering
	\begin{tikzpicture}[every node/.style={draw, thick, rounded corners, fill=blue!10, text width=2cm, align=center, minimum height=0.8cm}]
		\node (ch1) {\hyperlink{chapter 1}{فصل 1: مقدمه}};
		\node (ch2) [right=of ch1] {\hyperlink{chapter 2}{فصل 2: فرمول‌ها}};
		\node (ch3) [below=of ch1] {\hyperlink{chapter 3}{فصل 3: الگوریتم}};
		\node (ch4) [right=of ch3] {\hyperlink{chapter 4}{فصل 4: نتیجه}};
		\draw[->, thick] (ch1) -- (ch2);
		\draw[->, thick] (ch1) -- (ch3);
		\draw[->, thick] (ch2) -- (ch4);
		\draw[->, thick] (ch3) -- (ch4);
	\end{tikzpicture}
	\caption{گراف تعاملی فصول پروژه (با قابلیت کلیک)}
	\label{fig:graph-chapters}
\end{figure}

\section{پیشنهادات برای کارهای آینده}
با توجه به محدودیت‌های این پروژه، موارد زیر می‌تواند موضوع پژوهش‌های بعدی قرار گیرد:
\begin{enumerate}
	\item پیاده‌سازی نسخه‌های تصادفی گرادیان کاهشی (SGD) با مینی‌بچ و بررسی اثر نویز بر همگرایی.
	\item مقایسه روش گرادیان کاهشی با روش‌های مرتبه دوم مانند نیوتن و شبه‌نیوتن از نظر تعداد تکرار و زمان محاسبات.
	\item استفاده از روش‌های نرخ یادگیری تطبیقی (مانند Adam، RMSprop) برای مسائل با مقیاس‌های مختلف.
	\item اعمال الگوریتم بر روی مسائل بهینه‌سازی مقید با استفاده از روش‌های لاگرانژ و تابع جریمه.
\end{enumerate}

\section{جمع‌بندی نهایی}
در این پروژه کلیه اجزای درخواستی شامل شبه‌کد الگوریتم (\cref{alg:gd})، شکل با دو زیرشکل (\cref{fig:subfigures})، فلوچارت (\cref{fig:gd-flowchart})، جدول رنگی (\cref{tab:optim-compare})، سه فرمول (\cref{eq:grad-descent-long}، \cref{eq:piecewise-penalty} و \cref{eq:norm2}) و گراف تعاملی (\cref{fig:graph-chapters}) ارائه شد. همچنین با استفاده از دو مرجع \cite{boyd2004convex} و \cite{nocedal2006numerical} به منابع معتبر علمی ارجاع داده شده است. انتظار می‌رود این گزارش بتواند زمینه‌ساز مطالعه‌ی بیشتر در حوزه الگوریتم‌های گرادیان مبنا قرار گیرد.